\section{Gaia DR3 数据概述}

欧洲空间局的 Gaia 任务自 2013 年发射以来，
持续对全天空恒星源进行高精度天文测量，
其核心目标是建立一个与国际天球参考系（ICRF）严格对齐的高精度天文参考框架。
Gaia 观测数据以阶段性数据发布（Data Release）的形式向公众开放，
各次数据发布在观测时间跨度、数据处理流程以及可用天体参数方面均存在显著差异。

Gaia 的首次数据发布（DR1）主要提供了基于有限观测弧段的恒星位置，
其精度和参数完整性仍受到较大限制。
随着观测时间基线的延长和数据处理算法的持续改进，
第二次数据发布（DR2）首次为大规模恒星样本提供了完整的五参数天体测量解，
包括天球位置、自行和视差，
标志着 Gaia 星表在天文测量精度和参考框架稳定性方面的重大跃升。
基于 Gaia DR2，
Eggl 等人系统地更新了历史星表的去偏修正方案，
并证明以 Gaia 作为参考能够显著降低小天体光学观测中的区域性系统误差。

随后发布的早期第三次数据发布（EDR3）在观测时间跨度和系统误差控制方面进一步改进，
显著提升了恒星位置与自行的一致性和内部精度。
在此基础上，Gaia 第三次数据发布（DR3）
不仅继承了 EDR3 在天体测量精度上的改进成果，
还在数据处理一致性、系统误差建模以及附加天体物理参数方面进行了扩展，
从而形成了一套在精度、稳定性和覆盖范围上均优于先前版本的星表产品。

鉴于 Gaia DR3 在天体测量精度、系统误差控制以及参考框架稳定性方面的综合优势，
本工作选用 Gaia DR3 作为统一的参考星表，
用于评估和修正各类历史星表中的系统性位置与自行偏差。
通过将历史星表统一映射至 Gaia DR3 参考框架下，
可以在保持与既有 Gaia 时代去偏方法一致性的同时，
充分利用最新一代 Gaia 数据所提供的精度优势。

\newpage